% =====================================================================
%  qurra.tex — أوامر أسماء القراء والرواة والرموز
%  الاستعمال في النص:  \Asim{} و\Kisai{}  ← يظهر الاسم بلون صاحبه
% =====================================================================

% اسم ملوَّن عام:  \qname{لون}{الاسم}
\newcommand{\qname}[2]{\textcolor{c#1}{\textbf{#2}}}

% ---------- نافع وراوياه ----------
\newcommand{\Nafi}{\qname{Nafi}{نافع}}
\newcommand{\Qalun}{\qname{Nafi}{قالون}}
\newcommand{\Warsh}{\qname{Nafi}{ورش}}
% ---------- ابن كثير وراوياه ----------
\newcommand{\IbnKathir}{\qname{IbnKathir}{ابن كثير}}
\newcommand{\Bazzi}{\qname{IbnKathir}{البزي}}
\newcommand{\Qunbul}{\qname{IbnKathir}{قنبل}}
% ---------- أبو عمرو وراوياه ----------
\newcommand{\AbuAmr}{\qname{AbuAmr}{أبو عمرو}}
\newcommand{\AbiAmr}{\qname{AbuAmr}{أبي عمرو}}
\newcommand{\DuriAbuAmr}{\qname{AbuAmr}{الدوري عن أبي عمرو}}
\newcommand{\Susi}{\qname{AbuAmr}{السوسي}}
% ---------- ابن عامر وراوياه ----------
\newcommand{\IbnAmir}{\qname{IbnAmir}{ابن عامر}}
\newcommand{\Hisham}{\qname{IbnAmir}{هشام}}
\newcommand{\IbnDhakwan}{\qname{IbnAmir}{ابن ذكوان}}
% ---------- عاصم وراوياه ----------
\newcommand{\Asim}{\qname{Asim}{عاصم}}
\newcommand{\Shuba}{\qname{Asim}{شعبة}}
\newcommand{\Hafs}{\qname{Asim}{حفص}}
% ---------- حمزة وراوياه ----------
\newcommand{\Hamza}{\qname{Hamza}{حمزة}}
\newcommand{\Khalaf}{\qname{Hamza}{خلف}}
\newcommand{\Khallad}{\qname{Hamza}{خلاد}}
% ---------- الكسائي وراوياه ----------
\newcommand{\Kisai}{\qname{Kisai}{الكسائي}}
\newcommand{\AbulHarith}{\qname{Kisai}{أبو الحارث}}
\newcommand{\DuriKisai}{\qname{Kisai}{الدوري عن الكسائي}}

% ---------- الباقون ----------
\newcommand{\Baqun}{\textbf{\textcolor{black!70}{الباقون}}}

% ---------- الرموز داخل البيت ----------
% كلمة فيها رمز قارئ:  \rz{Kisai}{راويه}
\newcommand{\rz}[2]{\textcolor{c#1}{\textbf{#2}}}
% كلمة فيها رمز جماعة (ث خ ذ ظ غ ش أو صحبة، عم ...):  \rzg{صحبة}
\newcommand{\rzg}[1]{\textcolor{cGroup}{\textbf{#1}}}
% الكلمة القرآنية المختلف فيها داخل البيت:  \qw{مالك}
\newcommand{\qw}[1]{{\setlength{\fboxsep}{1.2pt}\colorbox{wordbg}{#1}}}

% ---------- الآيات ----------
% \ayah{مَالِكِ يَوْمِ الدِّينِ}  ← ﴿ ... ﴾ بخط المصحف
\newcommand{\ayah}[1]{{\quranfont\textcolor{ayahbracket}{﴿}#1\textcolor{ayahbracket}{﴾}}}
% \src{الفاتحة: ٤}  ← مرجع الآية
\newcommand{\src}[1]{{\small\textcolor{black!60}{[#1]}}}

% حرف رمز داخل الجداول بخلفية ملوّنة:  \rzl{Nafi}{أ}
\newcommand{\rzl}[2]{{\setlength{\fboxsep}{2pt}\colorbox{c#1!15}{\textcolor{c#1}{\textbf{\large #2}}}}}
\newcommand{\rzlg}[1]{{\setlength{\fboxsep}{2pt}\colorbox{cGroup!15}{\textcolor{cGroup}{\textbf{\large #1}}}}}

% سهم «يؤدي إلى» باتجاه القراءة العربية
\newcommand{\ila}{\ $\leftarrow$\ }
